\chapter{Introduction}
\label{chapter:introduction}

\textbf{Establishing of context: What is the HC?}

The Hadley cell (HC) is the dominant meridional circulation in the tropics. The region at the equator is associated with rising motion and collocated with inter-tropical convergence zone (ITCZ). The dry regions of subsidence in the subtropics are tied to the descending branch of the HC. In the annual mean there are two cell, which extend from \SIrange{0}{30}{\degree\N} in the Northern hemisphere and equivalently \SIrange{0}{30}{\degree S} in the Southern hemisphere. Throughout the annual cycle there are strong variations in the HC's extent and strength: As the maximum of insolation moves poleward, so does the ITCZ and the HCs, resulting in a strong winter and a weaker summer HC. This seasonal migration is related to the global monsoons.\\ 
The HC is responsible for the poleward transportation of angular momentum (and energy) and the equatorward transport of moisture. Therefore, the HC has strong implications for life on this planet. Its extent crucially modulates gloabl precipitation patterns \parencite{}. [Precipitation changes in the Sahel in the historical period. Paleoclimate HC and precipitation changes.] (this is unequal to the significance)The strength sets the intensity of updrafts (amount of precipitation), the moisture transport\\
Climate changes in the past had a significant influence on the HC altering its strength and extent. This is why there has been a large effort to find out the implications of human driven greenhouse gas emissions in the current era.


\textbf{Stating the Problem:}
To date, there has not been sufficient research that examines the connection between...

inter-model uncertainty
\CO is very uniformly distributed in the atmosphere. The forcing that it imposes is not. The HC strength response depends on meridional forcing structure. This has been shown for quadrupling of \CO only in certain latitude belts \parencite{kimWeakHadleyCell2022}. A forcing in midlatitudes weakens the Hadley circulation whereas a forcing at the equator leads to it strengthening. Thus, we want to investigate how \CO forcing in the vertical alters atmospheric circulation. We achieve this by quadrupling \CO only in the troposphere and only in the stratosphere using a pre-industrial control tropopause height. This approach is completely new and has net yet been described by the literature. Previous studies emphasise the importance of drivers related to the troposphere to explain the HC weakening \parencite{}. However, historical simulations show that also stratospheric forcing influences the general circulation. In partical the ozone hole in the Antarctic has lead to a strengthening of the circulation \parencite{}.  
We find that the HC response to this forcing is not what one would expect from the literature focusing tropospheric thermodynamical drivers for HC strength. Our studies disentangles the contributions of tropospheric and stratospheric \CO to the established HC weakening with uniform \CO distribution. and want to employ different frameworks for the HC dynamics. 

\textbf{Scope}
This study explores whether the simulated changes in HC strength can be explained using already established theories of HC dynamics.

\textbf{Significance}

% advance knowledge in the applicable field; that is, it revises or creates new knowledge (for example, the results will extend what is known about the applicability of a theory, the results are widely generalizable, or, for qualitative studies, transferable to other contexts)
We contribute to the advancement of the field by revising the applicability of different HC frameworks and simplified models. We transfer their results to the context of different vertical forcing structures
- foundational research that advances our physical understanding on the drivers of general circulation changes
\textbf{Overview over the study}
Is this really necessary?